% ECONOMICS_PROBLEM: {"id": "C78-237", "title": "Stable-roommates robustness to preference perturbations", "jel_code": "C78", "source_id": "OP-0561"}
% FIRST_POSTED: 2026-10-10
% ATTEMPT_STATUS: solved
% ENTRY_KIND: research_attempt
\documentclass[11pt]{article}
\usepackage[a4paper,margin=25mm]{geometry}
\usepackage[T1]{fontenc}
\usepackage{lmodern,microtype,amsmath,amssymb,amsthm}
\usepackage[hidelinks]{hyperref}
\newtheorem{theorem}{Theorem}
\newtheorem{lemma}{Lemma}
\newcommand{\NP}{\mathsf{NP}}
\newcommand{\FP}{\mathsf{FP}}
\title{Robust stable-roommates matchings:\\existence and optimization}
\author{Ji Ho Bae}
\date{10 October 2026}
\hypersetup{pdftitle={Robust stable-roommates matchings: existence and optimization},pdfauthor={Ji Ho Bae}}
\begin{document}
\maketitle

\paragraph{Authorship and provenance.}
Author: Ji Ho Bae. Prepared with GPT Codex assistance; the precise serving
revision was not exposed. The complete original argument is retained in
the \href{https://github.com/jbaelaw/robust-roommates-proof/blob/07e1c1655e18eaee6bb20468771afe6beabd93e5/paper/Proof.tex}{immutable TeX source} and
\href{https://github.com/jbaelaw/robust-roommates-proof/blob/07e1c1655e18eaee6bb20468771afe6beabd93e5/paper/Proof.pdf}{proof PDF}. This submission adds catalogue metadata,
provenance and a completion statement to that argument. Its preparation
included exact-target and primary-source checks and separate mathematical
reviews. This is a complete proof claim submitted for independent expert
review; community acceptance remains pending.

\begin{abstract}
For complete strict roommates preferences, a matching stable against
at most $d$ adjacent swaps in total across all lists can be found in polynomial time,
or its nonexistence certified. The algorithm adds symmetry constraints
to the rotation representation for robust marriage. Minimum-cost robust
matching is NP-hard at every fixed swap radius. Maximizing stability
probability under an explicit joint perturbation law is NP-hard even
when that law is a point mass. We give the corresponding verification,
decision and exact-search bounds.
\end{abstract}

\section{Model and classification}
Let $I$ give strict complete rankings on an even set of $N$ agents.
A matching partitions the agents into pairs. An unmatched pair blocks
when both agents prefer each other to their assigned partners.
A matching is $d$-robust if it is stable after every sequence of at most
$d$ adjacent swaps across all lists, where $d\geq0$ is encoded in binary.
For the separate probability question, specify a finite joint law
$\nu=\sum_{\ell=1}^s w_\ell\delta_{J_\ell}$ by listing complete strict
profiles and nonnegative binary rational weights summing to one. Write
\[
 r_\nu(M)=\sum_{\ell=1}^s w_\ell
       \mathbf1\{M\text{ is stable in }J_\ell\},\qquad M\in\mathcal S(I).
\]
Here $\mathcal S(I)$ is the set of originally stable matchings. The law
need not be supported in the swap ball; every listed profile is obtainable
by some finite sequence of adjacent swaps. Edge costs and objective
thresholds are binary rationals.

\begin{theorem}\label{classification}
For the model above, the following statements hold.
\begin{enumerate}
\item Verification, existence and construction of a $d$-robust matching
      are polynomial-time problems.
\item Deciding whether a $d$-robust matching of cost at most a given
      threshold exists is NP-complete for every fixed $d\geq0$, even
      with zero-one edge costs. It is also NP-complete for input $d$.
\item Evaluating $r_\nu(M)$ is polynomial-time. Deciding whether some
      $M\in\mathcal S(I)$ satisfies $r_\nu(M)\geq\lambda$ is NP-complete,
      even for a point-mass law and $\lambda=1$. Its maximization admits
      no deterministic polynomial-time multiplicative approximation
      with any finite factor unless $\mathsf P=\NP$.
\item Minimum cost subject to $M\in\mathcal S(I)$ and
      $r_\nu(M)\geq\lambda$ also has an NP-complete threshold problem.
      Both cost optimizations and probability maximization admit
      $\FP^{\NP}$ algorithms, reporting infeasibility when necessary.
\end{enumerate}
\end{theorem}

These statements answer catalogue problem C78-237 with the perturbation
representation specified above. They impose no restriction on the
original complete strict profile. The survey motivating the question
distinguishes preference perturbations from pair-failure robustness
\cite[\S3.3.2]{survey}; our model concerns the former.

\section{Verification and a polynomial existence criterion}
Number positions in agent $i$'s original list from one, writing $r_i(j)$.
For an unmatched pair put
\[
 b_M(i,j)=\max\{r_i(j)-r_i(M(i)),0\}
          +\max\{r_j(i)-r_j(M(j)),0\}.
\]
This is exactly the minimum total number of adjacent swaps needed to
make the pair block. If a name is already above the assigned partner,
its contribution is zero. Otherwise reversing their order costs their
position distance: they must cross each other, and each intervening
name must cross one of them. Moving the candidate directly above the
partner attains this bound. The two lists are distinct, so the costs add.
Consequently
\begin{equation}\label{radius}
 M\text{ is }d\text{-robust}
 \quad\Longleftrightarrow\quad
 b_M(i,j)>d\text{ for every unmatched pair }\{i,j\}.
\end{equation}
This gives a quadratic verification algorithm. For $N=0$ or $2$, the
unique matching is robust at every radius. For $N\geq4$, every pair cost
is at most $2(N-2)$, so $d\geq2(N-2)$ is immediately infeasible.

We use the following established marriage representation. For a strict
bipartite profile, allowing incomplete lists, stable matchings correspond
to predecessor-closed sets of rotations. With a variable $x_\rho$ saying
that rotation $\rho$ is eliminated, predecessor closure consists of
implications. Chen, Skowron and Sorge give polynomially many further
implications and unit clauses whose satisfying closed sets correspond
\emph{exactly} to all $d$-robust stable marriages
\cite[Proposition 3.9, Lemmas 3.15 and 3.19, and proof of Theorem 3.24]{css}.
In their notation, a relevant stable quadruple has associated rotations
$\pi,\rho$. When both exist its clause is $x_\pi\Rightarrow x_\rho$;
when only $\pi$ exists it is $\neg x_\pi$; when only $\rho$ exists it is
$x_\rho$; when neither exists the instance is infeasible. Only mutually
acceptable prospective blocking pairs are considered. Thus these
constraints form a polynomially constructible 2-CNF formula.

Replace each original agent $i$ by a left copy $L_i$ and a right copy
$R_i$. Both copies rank opposite-side names in $i$'s original order;
$L_iR_i$ is unacceptable. Run deferred acceptance on this doubled
marriage instance. If its matching is not perfect, reject: all stable
marriages match the same subsets of agents, whereas a stable roommate
matching would give a perfect stable marriage
\cite[Proposition 2.3]{css}. Otherwise every stable marriage is perfect.

Write such a matching as a permutation $p$, pairing $L_i$ with $R_{p(i)}$.
The partner name of $R_i$ is $p^{-1}(i)$. For each $i$ and
$t\in\{1,\ldots,N-1\}$ consider the two predicates
\[
 H^L_{i,t}=[r_i(p(i))\leq t],\qquad
 H^R_{i,t}=[r_i(p^{-1}(i))\leq t].
\]
Each predicate is constant, a rotation variable, or its negation. Left partners worsen as
rotations are eliminated; right partners improve. The rotations affecting
an agent are ordered and traverse her stable partners without repetition.
A nonconstant threshold is therefore crossed by exactly one rotation:
$H^L_{i,t}=\neg x_\alpha$ or $H^R_{i,t}=x_\beta$, respectively.
The crossing inequalities are $r_{\rm old}\leq t<r_{\rm new}$ on the
left and $r_{\rm old}>t\geq r_{\rm new}$ on the right.
The extreme stable matchings determine the constant cases, and the
explicit rotations determine the crossings
\cite[Proposition 3.6(ii),(iv)]{css}.

Add $H^L_{i,t}\leftrightarrow H^R_{i,t}$ for all $i,t$ to the marriage
robustness formula. Equality of two literals, including constants, is
2-CNF. Since the copies have identical strict rankings, these equalities
hold precisely when $p(i)=p^{-1}(i)$ for every $i$. Hence $p$ is an
involution; the forbidden diagonal excludes fixed points. Its pairs
define a roommate matching.

Ordinary stability is preserved by this correspondence. Moreover, for
a symmetric matching the cheapest attack by $L_iR_j$ has cost exactly
$b_M(i,j)$. Formula~\eqref{radius} therefore shows that robustness is
preserved too: each attack concerns only the two lists of its blocking pair.

The final 2-CNF formula is satisfiable if and only if a $d$-robust
roommate matching exists. A satisfying assignment specifies a closed
rotation set; eliminate it and project its symmetric matching.
Unsatisfiability certifies nonexistence. There are polynomially many
rotations, robustness clauses and threshold equalities. The preliminary
radius bound handles binary $d$ without enumerating perturbations.

\section{Minimum cost is hard at every fixed radius}
We first give the ordinary weighted-roommates reduction explicitly.
Its four-agent gadget is the Feder construction reproduced in
\cite[\S4.1, proof of Proposition 6]{adapting}.
For a Vertex Cover instance $G=(V,E)$, create $p_v,a_v,b_v,c_v$ for
each vertex $v$. Their preference prefixes are
\[
\begin{aligned}
 p_v&:a_v\succ \{p_u:uv\in E\}\succ c_v,&
 a_v&:b_v\succ p_v,\\
 b_v&:c_v\succ a_v,& c_v&:p_v\succ b_v.
\end{aligned}
\]
Order each neighbor block strictly and append all omitted names after
the displayed prefix. Thus every list is complete.

In any stable matching, $p_v$ must have partner $a_v$ or $c_v$.
Otherwise the pair $p_va_v$ forces $a_v$ to be paired with $b_v$, its
only preferred alternative; then $b_vc_v$ blocks. If $p_va_v$ is used,
$b_vc_v$ must also be used, since otherwise it blocks. If $p_vc_v$ is
used, $a_vb_v$ must be used for the same reason. The two possible modes
are therefore
\[
 A_v=\{p_va_v,b_vc_v\},\qquad B_v=\{p_vc_v,a_vb_v\}.
\]
Both are internally stable. A graph-edge pair $p_up_v$ blocks exactly
when both vertices use $B$. All other cross-gadget pairs are below
the assigned local partners. Thus stable matchings correspond exactly
to vertex covers $\{v:A_v\text{ is used}\}$. Give $p_va_v$ cost one
and every other edge cost zero. Matching cost equals cover size.

Fix any $d\geq0$, and let this core have $n\geq4$ agents. Choose an
even integer $m>\max\{d+1,(n-2)d\}$. Add dummies
$z_0,\ldots,z_{m-1}$, with indices modulo $m$. Dummy $z_a$ ranks its
antipodal mate $z_{a+m/2}$ first, then the other dummies by increasing
cyclic offset $1,\ldots,m-1$, omitting the mate, then all core agents.
In every core list insert $d$ distinct dummy names between each
consecutive pair of core names, and append unused dummies last.
This uses $(n-2)d<m$ dummies per list, with reuse across different lists.

Every stable matching pairs the mutually first-choice dummy mates.
It therefore restricts to a stable core matching, and every stable core
matching extends by the dummy pairs. Such an extension is $d$-robust.
For an unmatched core pair, stability provides an endpoint that prefers
its core partner; the inserted names make its attack cost at least $d+1$.
For a core--dummy pair the dummy's attack cost is at least $m-1>d$.
For an unmatched dummy pair of offset $h\ne m/2$, the forward cost is
$h$ if $h<m/2$, and $h-1$ otherwise. Adding the reverse cost gives
$m-1>d$. These cases exhaust all unmatched pairs, so
\eqref{radius} applies. Conversely, a robust matching is stable and
projects to the core.

Preserve core costs and give all dummy edges cost zero. The correspondence
preserves the objective. Its size is polynomial for each fixed $d$,
including zero. Vertex Cover hardness and the polynomial verifier prove
part~2. This also proves hardness for input $d$; the padding itself need
not be polynomial in an unbounded binary value of $d$.

\section{Stability probability and exact search}
Checking stability in all explicitly listed profiles evaluates $r_\nu(M)$
in polynomial bit time. For hardness use the NP-completeness of finding
a marriage matching stable in two complete strict profiles $P,Q$
\cite[Corollary 5]{joint}. Turn both into roommates profiles by putting
every opposite-side name above every same-side name. Any stable roommate
matching is bipartite: a same-side pair would force a same-side pair on
the other side, and one member of each pair would mutually prefer to
match together. Bipartite stability is precisely marriage stability.

Take $I$ to be the converted $P$, and let $\nu$ be a point mass at the
converted $Q$. Originally stable matchings exist because complete
marriage instances have stable matchings. Among them $r_\nu$ is either
zero or one, and its maximum is one exactly when $P,Q$ have a common
stable matching. This proves part~3, including the approximation claim:
any finite multiplicative guarantee on a positive optimum must return
value one and would decide the NP-complete problem. Giving every edge
cost zero also proves hardness with both a probability threshold and a
cost bound. All threshold problems are in NP by the explicit verifiers.

Finally, scale rational costs, or listed probabilities, by the product
of their positive denominators. This product has polynomial bit length.
Each objective is then integer-valued in a range with polynomially many
bits. Test feasibility first, then use binary search with NP queries to
find its exact optimum. Recover an optimizer by fixing matching edges
successively while retaining an optimal feasible completion; these
queries remain in NP. This proves the stated $\FP^{\NP}$ upper bounds.
These are polynomial-time Turing upper bounds; the decision hardness
reductions above are polynomial-time many-one reductions.

\section*{Completion Estimate}
100\%
The stated result resolves the original catalogue question under its full
model and quantifiers. This is the author's complete proof claim, pending
independent expert review.

\begin{thebibliography}{9}
\bibitem{survey} K. Cechl\'arov\'a, \'A. Cseh and D. F. Manlove.
Selected open problems in Matching Under Preferences.
\emph{Bulletin of the EATCS} 128 (2019), 14--38, \S3.3.2.
\href{https://hpi.de/friedrich/docs/publications/2019/BEATCS.pdf}{Primary manuscript}.
\bibitem{css} J. Chen, P. Skowron and M. Sorge.
Matchings under Preferences: Strength of Stability and Trade-Offs.
\emph{EC 2019}, 41--59.
\href{https://arxiv.org/pdf/1902.10535v2}{arXiv:1902.10535v2}.
The proposition and lemma numbers cited here refer to this full version.
\bibitem{adapting} R. Bredereck, J. Chen, D. Knop, J. Luo and R. Niedermeier.
Adapting Stable Matchings to Evolving Preferences.
\emph{AAAI 2020}, 1830--1837.
\href{https://arxiv.org/pdf/1907.01375v2}{arXiv:1907.01375v2}, \S4.1,
proof of Proposition 6.
\bibitem{joint} S. Miyazaki and K. Okamoto.
Jointly Stable Matchings.
\emph{ISAAC 2017}, LIPIcs 92, 56:1--56:12.
\href{https://doi.org/10.4230/LIPIcs.ISAAC.2017.56}{doi:10.4230/LIPIcs.ISAAC.2017.56}.
\end{thebibliography}
\end{document}
